\documentclass[nopdfbookmarks,a4paper, 11pt, twoside]{article}
\usepackage{multicol}
\usepackage{dsekcommon}
\usepackage{dsekdokument}
\usepackage{tabularx}
\usepackage[T1]{fontenc}
\usepackage[utf8]{inputenc}
\usepackage[swedish]{babel}
\usepackage{url}
\usepackage{xcolor}
\usepackage[dvipsnames]{xcolor}

\setheader{Strategiska målen}{Strategiska målen}{}

\title{Rasmus cookar caféet}
\author{Rassesven}

\begin{document}

\section*{Strategiska målen}

\section{Formalia}
\subsection{Sammanfattning}
Strategisk plan för D-sektionen med långsiktiga mål som sektionen skall jobba mot.

\subsection{Syfte}
Syftet med den strategiska planen är att hålla reda på långsiktiga projekt och att användas i framtagandet av verksamhetsplanen.

\subsection{Historik}
Utkast färdigställt av: Framtidsutskottet 2021
Uppdaterad på HTM-1 2023 av: Framtidsutskottet 2023
Uppdaterad på HTM-1 2024 av: Framtidsutskottet 2024
Uppdaterad på HTM-Val 2025 av: Framtidsutskottet 2025
Uppdaterad på HTM1 2026 av: Framtidsutskottet 2026

\subsection{Upplägg}
Målen i detta dokument är saker som sektionen vill jobba med under längre tid. Varje mål sträcker sig över tre åt och efter tre år kan det antingen anses uppnått eller skrivas tillbaka på en ny treårs-period.

Alla mål skall vara ordentligt motiverade samt ha en tydlig målsättning. Målen i den Strategiskaplanen ligger sedan till grund för en del av verksamhetsplanen.

\subsection{Utvärdering av målen}
Utvärdering av målen skall ske varje år. Ifall man i utvärderingen anser att målet är uppnått ska det lyftas på nästkommande sektionsmöte att plocka bort målet. I annat fall förs utvärderingen vidare som information till nästkommande år samt presenteras till sektionen.

\subsection{Ansvar}
Ansvariga för den Strategiska planen är Framtidsutskottet. Strategiska planen ägs av sektionsmötet.

\subsection{Nya mål}
Sektionsmötet beslutar på HTM1 varje år att antingen förlänga eller tillägga ett till två mål. Detta för att det inte överskrider att vara sex mål åt gången. Ansvariga att föreslå mål att skriva in är Framtidsutskottet.

\newpage

\section{Områden}
Nedan är alla mål som finns samt vilka år de är för.

\subsection{Studiebevakning 2022-2028}

De flesta av sektionens medlemmar studerar i minst 5 år. För att göra dessa åren så bra och nyttiga som möjligt är det viktigt att vi påverkar vår utbildning till det bättre. Till vår hjälp har vi Studierådet, som är sektionens viktigaste utskott, med studiebevakning som huvudansvar.

Även fast sektionens huvudsyfte är studiebevakning och sektionen inte kan existera utan studierådet är intresset för utskottet generellt ganska lågt. Även intresset för våra påverkansmetoder är lågt. Det ligger i sektionens medlemmars bästa att detta ändras på så att vår utbildning blir så bra och konkurrenskraftig som möjligt.

\begin{itemize}
    \item Öka intresset och synligheten för Studierådet
    \item Öka svarsfrekvensen på CEQ-enkäterna.
    \item Öka studenternas vetskap om medel att påverka deras studier.
\end{itemize}

\subsection{Sektionskultur och trivsel 2022-2028}

Under 2020 skickades en enkät ut från skolan där sektionen fick vara med och bidra med frågor. I svaren på enkäten framkom det att många C:are \&\& D:are tyckte att det var svårt att börja engagera sig i sektionen och därmed ta del av gemenskapen. Många tyckte även att iDét kändes ovälkomnande då alla som brukar vara där redan är goda vänner. Utöver detta har D-sektionen en generell kultur som främjats av högre given social status kring särskilda utskott, stort fokus på alkohol samt grabbighet. För att så många ska trivas så mycket som möjligt bör sektionen jobba för att ha en mer välkomnande och mindre grabbig kultur. I en annan enkät framkom det att D-sektionens medlemmar var de mest ensamma på hela LTH. Detta är såklart inte bra och är något som sektionen borde bli bättre på att motverka. Detta gäller alla medlemmar oavsett nivå av engagemang i sektionen.

Under våren 2023 genomfördes en undersökning, kallad D-Compassen, där trivseln och kultur på sektionen stod i fokus. Den gav en god insikt i hur medlemmar upplever sektionen utifrån de nämnda ämnena och är en viktig del att behålla i framtiden för att kontinuerligt undersöka trivseln och välmåendet hos medlemmarna på sektionen. Under våren 2023 genomfördes en undersökning, kallad D-Compassen, där trivseln och kultur på sektionen stod i fokus. Den gav en god insikt i hur medlemmar upplever sektionen utifrån de nämnda ämnena och är en viktig del att behålla i framtiden för att kontinuerligt undersöka trivseln och välmåendet hos medlemmarna på sektionen.

\begin{itemize}
    \item Så många av sektionens medlemmar som möjligt skall känna sig välkomna in i sektionen.
    \item Sektionen skall ha en sund alkoholkultur.
    \item Minska den grabbiga kulturen.
    \item Utjämna den givna sociala statusen mellan olika utskott.
    \item Minskad ensamhet bland D-sektionens medlemmar.
\end{itemize}

\subsection{Hållbar tillväxt 2024-2026}
År efter år kommer nya, ibland fler än dessförinnan, studenter till Lunds Universitet, LTH och framförallt till D-sektionen. För att kunna tillgodose fler medlemmars behov av gemenskap och rekreation bör sektionen växa på ett hållbart och produktivt sätt. Ett sätt att öka sektionens tillväxt och ge fler en möjlighet att engagera sig på sektionen är att lägga till fler poster. Dessa kan ha ett mer väldefinierat men begränsat ansvarsområde, vilket kan leda till ett engagemang som inte behöver gå ut över medlemmens välmående. Genom att sänka tröskeln för sektionsengagemang ger vi fler studenter möjligheten att bli uppslukade i allt vad vi kallar D-sek och studentengagemang.

\begin{itemize}
    \item Möjliggöra för att fler skall kunna engagera sig.
    \item Tillgodose fler medlemmars intressen genom att utöka sektionens verksamhet, eller inkludera i existerande verksamhet.
    \item Främja hållbart engagemang med fokus på individens hälsa.
\end{itemize}

\subsection{Råsa studentikåsa! 2024-2028}
Sektionen har en lång historia med många traditioner. Vissa av dessa är fortfarande i bruk och vissa har fallit ur tiden. För att inte tappa vårt arv och vår historia kan det med gott samvete belysas och främjas på olika sätt, såsom genom att tydligare presentera sektionens historia för medlemmar. Sektionen ska även arbeta aktivt för att behålla de hållbara traditioner som finns.
\begin{itemize}
    \item Arbeta för arkivering, fortskridande och bevaring av äldre traditioner
    \item Fortsätta arkivera sektionens historia samt historiskt material, både fysiskt och digitalt, och presentera detta på ett sätt som är tillgängligt för sektionens medlemmar
    \item Skapa nya och hållbara traditioner
    \item Modernisera icke hållbara traditioner
\end{itemize}

\subsection{Relationship goals 2026-2028}
En anledning till att studera är att få möjlighet att lära känna nya människor. Tyvärr har sektionens rykte på LTH gått upp och ned, vilket påverkar medlemmarnas möjlighet att anknyta kontakt. För att fortsätta att förbättra vårt rykte ska vi samarbeta med andra sektioner och främja relationen mellan sektionerna och dess medlemmar.

\begin{itemize}
    \item Skapa möjligheter för medlemmar för att möta andra sektioners medlemmar.
    \item Skapa och bevara bra relationer med organisationer relevanta för sektionen.
\end{itemize}

\end{document}